% !TeX root = ../main.tex
% chapters/conclusion.tex

\chapter{Conclusion}
\label{ch:conclusion}

% ── Summary ───────────────────────────────────────────────────────
\section{Summary}
\label{sec:conc-summary}

This thesis set out to modernize the domain-theoretic foundations
laid by Benton, Kennedy, and Varming~\cite{BKV2009}, extend them
with enriched-categorical and quantum structure, and demonstrate
that $\CPO$-enrichment provides a formal bridge between classical
and quantum domain theory.
The result is a library of approximately 15{,}700 lines of
\Rocq{}~9.0 code across 35~source files, with zero uses of
\texttt{Admitted} in the core and a single explicitly declared axiom.

\paragraph{Classical domain theory and PCF adequacy.}
The classical layer reproduces every construction from
BKV~\cite{BKV2009}---$\omega$-CPOs, products, sums, the lift monad,
function spaces, and the Kleene fixed-point theorem---on a
modernized Hierarchy Builder infrastructure
(\cref{sec:meth-hierarchy,sec:meth-constructions}).
The fixed-point operator is registered as a Scott-continuous map
$\fixp : [D \to_c D] \to_c D$, and we formalize the Scott topology
with an equivalence between computational and topological continuity
(\cref{sec:bg-scott-topology}).
Building on this infrastructure, we give a complete denotational
semantics for PCF$_v$ with intrinsically typed de~Bruijn syntax and
prove both soundness and computational adequacy via a logical
relation (\cref{sec:res-soundness,sec:res-adequacy}).
The adequacy proof establishes that a closed expression converges
operationally if and only if its denotation is defined---a
full convergence correspondence.

\paragraph{Enriched categories and domain equations.}
The enriched layer formalizes $\CPO$-enriched categories, locally
continuous functors, and enriched natural transformations with
vertical composition, whiskering, and pointwise suprema
(\cref{sec:meth-enriched,sec:meth-nat-trans}).
As a validation, we prove an enriched Yoneda lemma establishing a
$\CPO$ isomorphism between natural transformations out of a
representable functor and the representing object
(\cref{sec:res-yoneda}).
Building on the enriched layer, we formalize embedding-projection
pairs, mixed-variance locally continuous bifunctors, and the
approximation tower for recursive domain equations
(\cref{sec:meth-domain-equations}).
The bilimit specification packages a solution
$D_\infty \cong F(D_\infty, D_\infty)$ with \texttt{ROLL}/\texttt{UNROLL}
isomorphisms, isolating a single axiom (\texttt{bilimit\_exists})
whose consequences are fully derived
(\cref{sec:res-domain-equations}).

\paragraph{Quantum domain theory and the enrichment correspondence.}
Following Kornell, Lindenhovius, and Mislove~\cite{KLM2024}, we
formalize involutive quantales, quantum posets, quantum CPOs, and the
quantum Kleene fixed-point theorem
(\cref{sec:disc-quantum,sec:disc-qcpo}).
The thesis's central structural observation is that the hom-sets of
quantum Scott-continuous maps form classical $\omega$-CPOs under a
pointwise quantum order, making the category $\qCPO$
satisfy the axioms of a $\CPO$-enriched category---just as the
classical category $\CPO$ does
(\cref{sec:disc-quantum-enrichment}).
This means that locally continuous functors, enriched natural
transformations, and the machinery for recursive domain equations all
apply to quantum domains without modification.

\paragraph{The coinductive lift monad.}
As a complementary negative result, we develop the coinductive delay
monad in a self-contained 490-line file and prove that it cannot
carry a \texttt{CPO.type} structure under Leibniz equality
(\cref{sec:disc-lift-monad}).
The obstruction theorem (\texttt{delay\_CPO\_requires\_quotient}) is
fully proved and makes the design rationale for the flat lift
$D_\bot = \texttt{option}\;D$ precise and machine-checked.

% ── Future work ───────────────────────────────────────────────────
\section{Future Work}
\label{sec:conc-future}

Several concrete directions remain open, grouped below by
the layer of the library they concern.

\paragraph{Completing the classical layer.}
The most immediate task is proving \texttt{bilimit\_exists}: the
$\omega$-product CPO (a countable product of CPOs with componentwise
order and componentwise suprema) is the missing ingredient, and the
proof strategy is well understood (\cref{sec:disc-tradeoffs}).
A second direction is replacing \texttt{ClassicalEpsilon} in the
lift supremum construction with a decidability instance for finite
or discrete types, which would make those instances extractable to
OCaml or Haskell while leaving the general case classical.
On the engineering side, the two-tier architecture
(\cref{sec:disc-tradeoffs}) could be partially closed by registering
\texttt{cont\_fun} and \texttt{mono\_fun} as \HB{} structures, if
the single-carrier indexing constraint is relaxed in a future version
of Hierarchy Builder.
A related effort is a universe-polymorphism audit: the library
currently avoids \texttt{Set Universe Polymorphism} entirely,
foreclosing some generic constructions; enabling it would require
pervasive annotations across the hierarchy for what is currently
minimal practical gain, but would become worthwhile if the library
scales to larger categorical constructions.

\paragraph{Completing the quantum layer.}
The quantum layer has five open fronts.
First, the \HB{} instance registering $\qCPO$ as a
\texttt{CPOEnrichedCat} requires packaging the \texttt{q\_cont\_fun}
record as an \HB{}-compatible type; the mathematics is complete
(\cref{sec:disc-quantum-enrichment}).
Second, the library axiomatises the quantale~$Q$ abstractly but does
not construct the concrete example $Q = B(H)$; doing so would
require a Hilbert-space formalization not currently available in the
\Rocq{} ecosystem.
Third, Scott continuity for quantum maps---the quantum analogue of
the equivalence between computational and topological continuity
formalized in the classical layer---has not been developed.
Fourth, the classical lift $D_\bot = \texttt{option}\;D$ has no
direct quantum analogue; a quantum lift monad would need to interact
with the quantale-valued order in a way that remains an open
question.
Fifth, full recursive domain equations for quantum types would
require extending the bilimit machinery to the quantum setting,
using the $\CPO$-enrichment of $\qCPO$ to instantiate the existing
bilimit framework.

\paragraph{Language extensions.}
The stub files \texttt{QMiniCore\_Syntax.v} and
\texttt{QMiniCore\_Semantics.v} are reserved for a quantum
programming language whose types would be interpreted as quantum
CPOs.
A minimal viable language would include qubit allocation, unitary
gates, measurement, and a quantum-control branching construct
whose denotational semantics targets the quantum Kleene fixed-point
theorem.
Designing such a language and proving soundness and adequacy
analogous to the PCF$_v$ development would be the natural
culmination of the quantum layer.

% ── Closing remarks ───────────────────────────────────────────────
\section{Closing Remarks}
\label{sec:conc-closing}

The organizing principle of this thesis is that $\CPO$-enrichment is
a minimal and powerful abstraction: it captures exactly the structure
needed for fixed-point semantics, recursive domain equations, and
enriched categorical reasoning, and it applies uniformly to both
classical and quantum domains.
By building the enriched layer as a generic framework and then
demonstrating that both $\CPO$ and $\qCPO$ satisfy its axioms, the
thesis shows that a single mechanized infrastructure can serve as a
foundation for programming language semantics across the
classical--quantum divide.
The gaps that remain---a concrete Hilbert-space quantale, a quantum
lift monad, a quantum programming language---are gaps in
\emph{content}, not in \emph{architecture}: the framework is in
place, and filling these gaps requires extending the library rather
than redesigning it.
